\section{Projektinitiering}

% jonathan
% alt+z for line wrap :)
Gruppen valgte at holde problemfelt og problemformulering mere åbne for at undgå at definere den endelige løsning for tidligt i projektet. Dette var inspireret af PBL-tilgangen, hvor problemfeltet skal beskrive konteksten uden at skabe tunnelsyn omkring en bestemt løsning eller historisk retning.

Gruppen valgte efterfølgende at tage udgangspunkt i razziaen fra oktober 1943, da denne begivenhed repræsenterede en markant og dramatisk del af skolens historie under besættelsen. Begivenheden gav samtidig mulighed for at kombinere historisk formidling med multiplayer gameplay og samarbejdsmekanikker i Unreal Engine 5. Projektets historiske tema førte naturligt til flere etiske overvejelser. Gruppen diskuterede tidligt i projektet, hvordan historiske roller skulle repræsenteres i spillet. Selvom projektet var inspireret af virkelige hændelser fra besættelsen, var målet ikke at skabe en fuldstændig historisk simulation, men derimod en multiplayeroplevelse inspireret af skolens historie. Der opstod derfor refleksioner omkring balancen mellem historisk autenticitet og spildesign. Gruppen valgte blandt andet, at spillere ikke skulle fremstilles som elever i nazistiske roller. Denne beslutning blev taget for at undgå en uhensigtsmæssig identifikation mellem nuværende elever og historiske aktører fra besættelsen. Samtidig blev flere gameplay-elementer, såsom abilities og movement-mekanikker, designet ud fra etablerede multiplayergenrer fremfor historisk realisme. Dette gjorde det muligt at prioritere et engagerende gameplay-loop, samtidig med at spillet stadig bevarede en historisk inspiration.